\documentclass[11pt,a4paper]{article}

\usepackage[utf8]{inputenc}
\usepackage[T1]{fontenc}
\usepackage[provide=*,english]{babel}
\usepackage{lmodern}
\usepackage{microtype}
\usepackage{amsmath,amssymb,mathtools}
\usepackage{geometry}
\usepackage{xcolor}
\usepackage{enumitem}
\usepackage{booktabs}
\usepackage{fancyhdr}
\usepackage[hidelinks]{hyperref}

\geometry{left=2.3cm,right=2.3cm,top=2.2cm,bottom=2.25cm,headheight=14pt}
\setlength{\parindent}{0pt}
\setlength{\parskip}{0.38em}
\setlist[itemize]{topsep=0.25em,itemsep=0.18em,leftmargin=1.5em}
\setlist[enumerate]{topsep=0.25em,itemsep=0.22em,leftmargin=1.8em}
\allowdisplaybreaks

\definecolor{azul}{RGB}{34,73,110}
\definecolor{gris}{RGB}{90,96,102}

\pagestyle{fancy}
\fancyhf{}
\fancyhead[L]{\small\color{gris}GC2D: ABBA method}
\fancyhead[R]{\small\color{gris}Projection and constant Jacobian}
\fancyfoot[C]{\thepage}
\renewcommand{\headrulewidth}{0.3pt}

\newcommand{\R}{\mathbb{R}}
\newcommand{\Id}{I}
\newcommand{\Y}{\mathcal{Y}}
\newcommand{\M}{\mathcal{M}}
\newcommand{\Om}{\Omega_{\times}}
\newcommand{\PhiAB}{\Phi^{\mathrm{ABBA}}_{h,t_n}}
\newcommand{\norm}[1]{\left\lVert #1\right\rVert}

\begin{document}

\begin{center}
  {\LARGE\bfseries\color{azul}ABBA Method with Hairer's Symmetric Projection}\par
  \vspace{0.35em}
  {\large Adaptation to GC2D with $H_\mu=0$ and a constant simultaneous formulation}\par
  \vspace{0.45em}
  {\small Complete replacement of BM4 by an $A$--$B$--$B$--$A$ architecture}
\end{center}

\vspace{0.5em}

\section{Guiding-center system and duplicated phase space}

The physical state is
\[
 z=(X,Y)\in\R^2,
\]
and, with the convention used here, it satisfies
\[
 \dot X=-\partial_Y\phi(t,X,Y),
 \qquad
 \dot Y=\partial_X\phi(t,X,Y).
\]
We define
\[
 J_0=
 \begin{pmatrix}
  0&-1\\
  1&0
 \end{pmatrix},
 \qquad
 f(z,t)=J_0\nabla_z\phi(z,t),
\]
so that $\dot z=f(z,t)$. Its spatial Jacobian is
\[
 W(z,t):=D_zf(z,t)=J_0\nabla_z^2\phi(z,t).
\]
Since the Hessian of $\phi$ is symmetric,
\[
 W(z,t)^TJ_0+J_0W(z,t)=0. \tag{1}
\]

We duplicate the state:
\[
 \Y=(u,v)=(X_1,Y_1,X_2,Y_2)\in\R^4,
 \qquad u,v\in\R^2.
\]
The duplicated dynamics associated with the subflows used below is
\[
 \dot u=f(v,t),
 \qquad
 \dot v=f(u,t).
\]
If initially $u=v=z$, both equations coincide with the physical dynamics.

The physical state is identified with the diagonal manifold
\[
 \M=\{(u,v)\in\R^4:u=v\}.
\]
We write it as
\[
 g(\Y)=G\Y=u-v,
 \qquad
 G=(\Id_2\; -\Id_2),
 \qquad
 GG^T=2\Id_2.
\]
For $\mu=(\mu_X,\mu_Y)^T$,
\[
 G^T\mu=
 \begin{pmatrix}\mu\\-\mu\end{pmatrix}
 =(\mu_X,\mu_Y,-\mu_X,-\mu_Y)^T.
\]

The constraints are linear. Therefore,
\[
 \nabla^2g_1=\nabla^2g_2=0
\]
and, for every $\mu$,
\[
 \boxed{
 H_\mu(\Y)=\sum_{j=1}^2\mu_j\nabla^2g_j(\Y)=0.
 }
\]
This does not mean that $\mu=0$: $H_\mu$ is identically zero because $G=Dg$ is constant.

\section{The explicit submaps \texorpdfstring{$A$}{A} and \texorpdfstring{$B$}{B}}

For a substep length $s$ and an evaluation time $t$, we define
\[
 A_{s,t}(u,v)=\bigl(u+s f(v,t),v\bigr), \tag{2}
\]
\[
 B_{s,t}(u,v)=\bigl(u,v+s f(u,t)\bigr). \tag{3}
\]
The map $A$ advances only the first copy using the second one; $B$ advances only the
second copy using the first one. Both maps are explicit.

Their Jacobians are
\[
 DA_{s,t}(u,v)=
 \begin{pmatrix}
  \Id_2&sW(v,t)\\
  0&\Id_2
 \end{pmatrix},
 \qquad
 DB_{s,t}(u,v)=
 \begin{pmatrix}
  \Id_2&0\\
  sW(u,t)&\Id_2
 \end{pmatrix}. \tag{4}
\]

In the duplicated phase space, we use the skew-symmetric matrix
\[
 \Om=
 \begin{pmatrix}
  0&J_0\\
  J_0&0
 \end{pmatrix}.
\]
Identity (1) directly implies
\[
 (DA_{s,t})^T\Om DA_{s,t}=\Om,
 \qquad
 (DB_{s,t})^T\Om DB_{s,t}=\Om.
\]
Thus, each submap is symplectic with respect to $\Om$, and any composition of these maps is
also symplectic in the duplicated phase space.

Moreover, the variables used to evaluate the field remain fixed in each submap. Therefore,
\[
 \boxed{A_{s,t}^{-1}=A_{-s,t},\qquad B_{s,t}^{-1}=B_{-s,t}.} \tag{5}
\]

\section{The nonautonomous \texorpdfstring{$A$--$B$--$B$--$A$}{A-B-B-A} step}

Let
\[
 t_{n+1}=t_n+h,
 \qquad
 s=\frac h2.
\]
The base map that replaces BM4 is
\[
 \boxed{
 \PhiAB
 =A_{s,t_{n+1}}
  \circ B_{s,t_{n+1}}
  \circ B_{s,t_n}
  \circ A_{s,t_n}.
 } \tag{6}
\]
Compositions are read from right to left. Therefore, the effective order is
\[
 A_{s,t_n}\longrightarrow B_{s,t_n}
 \longrightarrow B_{s,t_{n+1}}\longrightarrow A_{s,t_{n+1}}.
\]

Starting from $(u^{(0)},v^{(0)})$, the four stages are
\begin{align}
 u^{(1)}
 &=u^{(0)}+s f\bigl(v^{(0)},t_n\bigr), &&(A_1) \tag{7a}\\
 v^{(1)}
 &=v^{(0)}+s f\bigl(u^{(1)},t_n\bigr), &&(B_1) \tag{7b}\\
 v^{(2)}
 &=v^{(1)}+s f\bigl(u^{(1)},t_{n+1}\bigr), &&(B_2) \tag{7c}\\
 u^{(2)}
 &=u^{(1)}+s f\bigl(v^{(2)},t_{n+1}\bigr). &&(A_2) \tag{7d}
\end{align}
The output is
\[
 \PhiAB(u^{(0)},v^{(0)})=(u^{(2)},v^{(2)}).
\]

The evaluations at the two time endpoints are essential in the nonautonomous problem.
If every stage were evaluated at $t_n$, time symmetry would be lost. In the autonomous
case, the two central stages can be combined, and one recovers
\[
 A_{h/2}\longrightarrow B_h\longrightarrow A_{h/2}.
\]

Since $A$ and $B$ are symplectic with respect to $\Om$, the complete map is also symplectic:
\[
 \boxed{
 (D\PhiAB)^T\Om D\PhiAB=\Om.
 } \tag{8}
\]

\section{Proof of the symmetry of ABBA}

A time-dependent method is symmetric if
\[
 \Phi^{\mathrm{ABBA}}_{-h,t_{n+1}}
 =\left(\Phi^{\mathrm{ABBA}}_{h,t_n}\right)^{-1}. \tag{9}
\]
Using (5) and reversing the order of composition in (6),
\begin{align*}
 \left(\Phi^{\mathrm{ABBA}}_{h,t_n}\right)^{-1}
 &={}
 A_{-s,t_n}
 \circ B_{-s,t_n}
 \circ B_{-s,t_{n+1}}
 \circ A_{-s,t_{n+1}}.
\end{align*}
We now start at $t_{n+1}$ with a step of size $-h$. Its final time is
$t_{n+1}-h=t_n$, and definition (6) gives precisely
\[
 \Phi^{\mathrm{ABBA}}_{-h,t_{n+1}}
 =A_{-s,t_n}
  \circ B_{-s,t_n}
  \circ B_{-s,t_{n+1}}
  \circ A_{-s,t_{n+1}}.
\]
Therefore,
\[
 \boxed{
 \Phi^{\mathrm{ABBA}}_{-h,t_{n+1}}
 =\left(\Phi^{\mathrm{ABBA}}_{h,t_n}\right)^{-1}.
 }
\]
The architecture is palindromic when both the sign of the step and the two endpoint times
are reversed simultaneously.

\section{Hairer's symmetric projection around ABBA}

Suppose that
\[
 \Y_n=(z_n,z_n)\in\M
\]
is known. For a multiplier $\mu\in\R^2$, we first correct the input:
\[
 \widehat{\Y}_n(\mu)
 =\Y_n+G^T\mu
 =(z_n+\mu,z_n-\mu).
\]
We write
\[
 u^{(0)}=z_n+\mu,
 \qquad
 v^{(0)}=z_n-\mu,
\]
and apply the four stages in (7). The auxiliary result is
\[
 \widetilde{\Y}_{n+1}(\mu)
 =\PhiAB\bigl(\Y_n+G^T\mu\bigr)
 =(u^{(2)}(\mu),v^{(2)}(\mu)).
\]
The same correction is then applied at the output:
\[
 \Y_{n+1}(\mu)
 =\widetilde{\Y}_{n+1}(\mu)+G^T\mu
 =\bigl(u^{(2)}+\mu,v^{(2)}-\mu\bigr). \tag{10}
\]
Imposing $\Y_{n+1}\in\M$ is equivalent to
\[
 u^{(2)}(\mu)+\mu=v^{(2)}(\mu)-\mu.
\]
Thus, the problem is reduced exactly to two equations:
\[
 \boxed{
 R(\mu)
 :=u^{(2)}(\mu)-v^{(2)}(\mu)+2\mu=0.
 } \tag{11}
\]
In matrix notation,
\[
 R(\mu)
 =G\PhiAB\bigl(\Y_n+G^T\mu\bigr)+2\mu. \tag{12}
\]
Once the root $\mu_*$ has been found, the new physical state is
\[
 \boxed{
 z_{n+1}=u^{(2)}(\mu_*)+\mu_*
          =v^{(2)}(\mu_*)-\mu_*.
 } \tag{13}
\]

\section{Exact Jacobian of the ABBA map}

For a particular evaluation of (7), we define
\begin{align*}
 W_1&=W\bigl(v^{(0)},t_n\bigr),
 &W_2&=W\bigl(u^{(1)},t_n\bigr),\\
 W_3&=W\bigl(u^{(1)},t_{n+1}\bigr),
 &W_4&=W\bigl(v^{(2)},t_{n+1}\bigr),
\end{align*}
and abbreviate
\[
 S:=W_2+W_3.
\]
The Jacobians of the four stages, evaluated at the points actually traversed, are
\[
 M_{A_1}=
 \begin{pmatrix}
  \Id_2&sW_1\\0&\Id_2
 \end{pmatrix},
 \qquad
 M_{B_1}=
 \begin{pmatrix}
  \Id_2&0\\sW_2&\Id_2
 \end{pmatrix},
\]
\[
 M_{B_2}=
 \begin{pmatrix}
  \Id_2&0\\sW_3&\Id_2
 \end{pmatrix},
 \qquad
 M_{A_2}=
 \begin{pmatrix}
  \Id_2&sW_4\\0&\Id_2
 \end{pmatrix}.
\]
By the chain rule,
\[
 \boxed{
 D\PhiAB=M_{A_2}M_{B_2}M_{B_1}M_{A_1}.
 } \tag{14}
\]
The two central matrices satisfy
\[
 M_{B_2}M_{B_1}
 =
 \begin{pmatrix}
  \Id_2&0\\sS&\Id_2
 \end{pmatrix}.
\]
Multiplying the four blocks gives the exact $4\times4$ Jacobian:
\[
 \boxed{
 D\PhiAB=
 \begin{pmatrix}
  \Id_2+s^2W_4S
  &s(W_1+W_4)+s^3W_4SW_1\\[2mm]
  sS
  &\Id_2+s^2SW_1
 \end{pmatrix}.
 } \tag{15}
\]
The Jacobian of the map has not been replaced by the identity: (15) contains exactly all
the dependencies generated by the four stages.

\section{Exact Jacobian of the reduced residual}

From
\[
 \frac{\partial(u^{(0)},v^{(0)})}{\partial\mu}
 =\begin{pmatrix}\Id_2\\-\Id_2\end{pmatrix}=G^T
\]
and from (12), we obtain
\[
 \boxed{
 R'(\mu)=G\,D\PhiAB\bigl(\Y_n+G^T\mu\bigr)G^T+2\Id_2.
 } \tag{16}
\]
If we write (15) as
\[
 D\PhiAB=
 \begin{pmatrix}A&B\\C&D\end{pmatrix},
\]
then
\[
 G\,D\PhiAB\,G^T=A-B-C+D.
\]
Substituting the four blocks in (15) yields
\[
 \boxed{
 \begin{aligned}
 R'(\mu)={}&4\Id_2
 -s\bigl(W_1+W_2+W_3+W_4\bigr)\\
 &+s^2\bigl[W_4(W_2+W_3)+(W_2+W_3)W_1\bigr]\\
 &-s^3W_4(W_2+W_3)W_1.
 \end{aligned}
 } \tag{17}
\]
Since $s=h/2$, this can also be written as
\[
 \boxed{
 \begin{aligned}
 R'(\mu)={}&4\Id_2
 -\frac h2\bigl(W_1+W_2+W_3+W_4\bigr)\\
 &+\frac{h^2}{4}
 \bigl[W_4(W_2+W_3)+(W_2+W_3)W_1\bigr]\\
 &-\frac{h^3}{8}W_4(W_2+W_3)W_1.
 \end{aligned}
 } \tag{18}
\]
The matrices $W_i$ need not commute: the order of the products in (17) and (18) must be
preserved exactly.

\section{Symmetry of the complete projected method}

The equation defining the projected step can be written as
\[
 \boxed{
 \Y_{n+1}-G^T\mu
 =\Phi^{\mathrm{ABBA}}_{h,t_n}
  \bigl(\Y_n+G^T\mu\bigr),
 \qquad
 \Y_n,\Y_{n+1}\in\M.
 } \tag{19}
\]
Applying the inverse of ABBA and using its symmetry,
\[
 \Y_n+G^T\mu
 =\Phi^{\mathrm{ABBA}}_{-h,t_{n+1}}
  \bigl(\Y_{n+1}-G^T\mu\bigr).
\]
We define the multiplier of the reverse step as
\[
 \mu_{\mathrm{inv}}=-\mu.
\]
Then
\[
 \Y_n-G^T\mu_{\mathrm{inv}}
 =\Phi^{\mathrm{ABBA}}_{-h,t_{n+1}}
  \bigl(\Y_{n+1}+G^T\mu_{\mathrm{inv}}\bigr),
\]
which is exactly equation (19) applied backward. Therefore, if the root is locally unique
and is solved exactly,
\[
 \boxed{
 \Psi_{-h,t_{n+1}}\circ\Psi_{h,t_n}=\Id.
 } \tag{20}
\]
Here $\Psi$ denotes ABBA together with Hairer's projection. In an implementation, the
iteration is stopped at a finite tolerance; reversibility is preserved up to that tolerance.

\section{Simultaneous formulation and constant approximation}

Alternatively, $\Y_{n+1}$ can be retained as an unknown. We define
\[
 d(\Y_{n+1},\mu)
 :=\Y_{n+1}-G^T\mu
   -\PhiAB\bigl(\Y_n+G^T\mu\bigr).
\]
We seek
\[
 d(\Y_{n+1},\mu)=0,
 \qquad
 g(\Y_{n+1})=0.
\]
Since $G$ is constant and $H_\mu=0$,
\[
 D_{\Y_{n+1}}d=\Id_4,
 \qquad
 D_\mu d
 =-\left[\Id_4+D\PhiAB(\widehat{\Y}_n)\right]G^T,
\]
\[
 D_{\Y_{n+1}}g=G,
 \qquad
 D_\mu g=0.
\]
The exact simultaneous Newton iteration solves
\[
 \boxed{
 \begin{pmatrix}
  \Id_4&-\left[\Id_4+D\PhiAB(\widehat{\Y}_n)\right]G^T\\[2pt]
  G&0
 \end{pmatrix}
 \begin{pmatrix}
  \Delta\Y_{n+1}\\
  \Delta\mu
 \end{pmatrix}
 =-
 \begin{pmatrix}
  d(\Y_{n+1},\mu)\\
  g(\Y_{n+1})
 \end{pmatrix}.
 } \tag{21}
\]
Its Schur complement is
\[
 G\left[\Id_4+D\PhiAB(\widehat{\Y}_n)\right]G^T
 =2\Id_2+GD\PhiAB(\widehat{\Y}_n)G^T
 =R'(\mu).
\]
Therefore, the reduced and simultaneous formulations use exactly the same Jacobian.

\end{document}
